% GENERATED FILE. Edit canonical lessons, then run npm run generate:books.
% canonical-source-hash: fnv1a64:15787a852b17ab88
% canonical-lessons: IT-C505-collezionare, IT-C505-collezione, IT-C505-puzzle, IT-C505-modellismo, IT-C505-assemblare

\chapter{Collezionisti}
\label{ch:it-collezionisti}
\hlchaptermodality{505}

\begin{chapteropening}
\textbf{By the end of this chapter:} \emph{I can talk about collecting and building models.}

Complete the last lesson of chapter 505: I can talk about collecting and building models.
\end{chapteropening}

\section[collezionare]{collezionare — to collect (as a hobby)}

\label{lesson:IT-C505-collezionare}



\begin{quote}
Before the new one: say the Italian for bet, then the Italian for crossword.
\end{quote}

\subsection*{You'll want to know: collezionare}
\textbf{collezionare} — \textquotedblleft{}to collect (as a hobby)\textquotedblright{}.

\begin{quote}
Da bambino collezionavo francobolli.
\end{quote}

\begin{quote}
\emph{As a child I collected stamps.}
\end{quote}

For things kept as a hobby; to pick up, say \emph{raccogliere}.

\begin{cousinweb}
\emph{collezionare} is built on \emph{collezione} \textquotedblleft{}collection\textquotedblright{}, a learned borrowing from Latin \emph{collēctiō} \textquotedblleft{}a gathering\textquotedblright{} (EtymDB), from \emph{colligere} \textquotedblleft{}to gather\textquotedblright{}.
\end{cousinweb}

\begin{grammarlens}[title={What you've built}]
One more word for telling, explaining and describing at B1.
\end{grammarlens}

\subsection*{Your turn}
\begin{itemize}
\raggedright
  \item \emph{Say it:} \emph{collezionare}
  \item \emph{Say it:} \emph{collezionare}, once more
  \item \emph{Say it:} a sentence of your own with \emph{collezionare}
  \item \emph{Recall:} say the Italian for bet, then the Italian for crossword, then say \emph{collezionare} again
\end{itemize}

\begin{tcolorbox}[breakable,colback=teal!4,colframe=teal!35!black,title={Before you move on}]
What does collezionare mean? (\textquotedblleft{}To collect (as a hobby)\textquotedblright{}.) Say it once more.
\end{tcolorbox}
\section[la collezione]{la collezione — collection}

\label{lesson:IT-C505-collezione}



\begin{quote}
Before the new one: say the Italian for crossword, then the Italian for to collect.
\end{quote}

\subsection*{You'll want to know: la collezione}
\textbf{la collezione} — \textquotedblleft{}collection\textquotedblright{}.

\begin{quote}
Ha una collezione di dischi in vinile.
\end{quote}

\begin{quote}
\emph{He has a collection of vinyl records.}
\end{quote}

\emph{Fare collezione di} — to collect.

\begin{cousinweb}
\emph{collezione} is a learned borrowing from Latin \emph{collēctiō} \textquotedblleft{}a gathering\textquotedblright{} (EtymDB), from \emph{colligere} \textquotedblleft{}to gather\textquotedblright{}, built on \emph{legere} \textquotedblleft{}to gather\textquotedblright{}.
\end{cousinweb}

\begin{grammarlens}[title={What you've built}]
One more word for telling, explaining and describing at B1.
\end{grammarlens}

\subsection*{Your turn}
\begin{itemize}
\raggedright
  \item \emph{Say it:} \emph{la collezione}
  \item \emph{Say it:} \emph{la collezione}, once more
  \item \emph{Say it:} a sentence of your own with \emph{la collezione}
  \item \emph{Recall:} say the Italian for crossword, then the Italian for to collect, then say \emph{la collezione} again
\end{itemize}

\begin{tcolorbox}[breakable,colback=teal!4,colframe=teal!35!black,title={Before you move on}]
What does la collezione mean? (\textquotedblleft{}Collection\textquotedblright{}.) Say it once more.
\end{tcolorbox}
\section[il puzzle]{il puzzle — jigsaw puzzle}

\label{lesson:IT-C505-puzzle}



\begin{quote}
Before the new one: say the Italian for to collect, then the Italian for collection.
\end{quote}

\subsection*{You'll want to know: il puzzle}
\textbf{il puzzle} — \textquotedblleft{}jigsaw puzzle\textquotedblright{}.

\begin{quote}
Abbiamo fatto un puzzle di mille pezzi.
\end{quote}

\begin{quote}
\emph{We did a thousand-piece jigsaw.}
\end{quote}

Said as in English; the plural is \emph{i puzzle}.

\begin{cousinweb}
\emph{puzzle} was borrowed from English \emph{puzzle} (EtymDB); the origin of the English word is not settled.
\end{cousinweb}

\begin{grammarlens}[title={What you've built}]
One more word for telling, explaining and describing at B1.
\end{grammarlens}

\subsection*{Your turn}
\begin{itemize}
\raggedright
  \item \emph{Say it:} \emph{il puzzle}
  \item \emph{Say it:} \emph{il puzzle}, once more
  \item \emph{Say it:} a sentence of your own with \emph{il puzzle}
  \item \emph{Recall:} say the Italian for to collect, then the Italian for collection, then say \emph{il puzzle} again
\end{itemize}

\begin{tcolorbox}[breakable,colback=teal!4,colframe=teal!35!black,title={Before you move on}]
What does il puzzle mean? (\textquotedblleft{}Jigsaw puzzle\textquotedblright{}.) Say it once more.
\end{tcolorbox}
\section[il modellismo]{il modellismo — model making}

\label{lesson:IT-C505-modellismo}



\begin{quote}
Before the new one: say the Italian for collection, then the Italian for jigsaw puzzle.
\end{quote}

\subsection*{You'll want to know: il modellismo}
\textbf{il modellismo} — \textquotedblleft{}model making\textquotedblright{}.

\begin{quote}
La sua passione è il modellismo: costruisce treni in miniatura.
\end{quote}

\begin{quote}
\emph{His passion is model making: he builds miniature trains.}
\end{quote}

The person is \emph{il modellista}.

\begin{cousinweb}
\emph{modellismo} is built on \emph{modello} \textquotedblleft{}model\textquotedblright{}, from Vulgar Latin \emph{modellus} (EtymDB), a little form of Latin \emph{modulus} \textquotedblleft{}a small measure\textquotedblright{} (Lewis \& Short).
\end{cousinweb}

\begin{grammarlens}[title={What you've built}]
One more word for telling, explaining and describing at B1.
\end{grammarlens}

\subsection*{Your turn}
\begin{itemize}
\raggedright
  \item \emph{Say it:} \emph{il modellismo}
  \item \emph{Say it:} \emph{il modellismo}, once more
  \item \emph{Say it:} a sentence of your own with \emph{il modellismo}
  \item \emph{Recall:} say the Italian for collection, then the Italian for jigsaw puzzle, then say \emph{il modellismo} again
\end{itemize}

\begin{tcolorbox}[breakable,colback=teal!4,colframe=teal!35!black,title={Before you move on}]
What does il modellismo mean? (\textquotedblleft{}Model making\textquotedblright{}.) Say it once more.
\end{tcolorbox}
\section[assemblare]{assemblare — to assemble, to put together}

\label{lesson:IT-C505-assemblare}



\begin{quote}
Before the new one: say the Italian for jigsaw puzzle, then the Italian for model making.
\end{quote}

\subsection*{You'll want to know: assemblare}
\textbf{assemblare} — \textquotedblleft{}to assemble, to put together\textquotedblright{}.

\begin{quote}
Ho assemblato il mobile da solo in due ore.
\end{quote}

\begin{quote}
\emph{I put the piece of furniture together by myself in two hours.}
\end{quote}

For furniture or parts; \emph{montare} means the same in speech.

\begin{cousinweb}
\emph{assemblare} came from French \emph{assembler} \textquotedblleft{}to bring together\textquotedblright{}, usually traced to Latin \emph{simul} \textquotedblleft{}together\textquotedblright{} (Lewis \& Short).
\end{cousinweb}

\begin{grammarlens}[title={What you've built}]
One more word for telling, explaining and describing at B1.
\end{grammarlens}

\subsection*{Your turn}
\begin{itemize}
\raggedright
  \item \emph{Say it:} \emph{assemblare}
  \item \emph{Say it:} \emph{assemblare}, once more
  \item \emph{Say it:} a sentence of your own with \emph{assemblare}
  \item \emph{Recall:} say the Italian for jigsaw puzzle, then the Italian for model making, then say \emph{assemblare} again
\end{itemize}

\begin{tcolorbox}[breakable,colback=teal!4,colframe=teal!35!black,title={Before you move on}]
What does assemblare mean? (\textquotedblleft{}To assemble, to put together\textquotedblright{}.) Say it once more.
\end{tcolorbox}
